{\Large \textbf{Abstract}}

\bigskip

Recreational and competitive athletes rely on several apps to track training load, monitor recovery, plan routes, and schedule sessions. This data stays split across proprietary ecosystems (Strava for activities, Garmin for wellness, weather services for conditions, calendars for availability), so athletes have to cross-reference them by hand, which is slow and error-prone. This paper presents \textit{Training Copilot}, an open-source, multi-agent sports analytics platform that brings these sources together behind one conversational interface. Training Copilot uses the Model Context Protocol (MCP) as a uniform integration layer: each external API is wrapped in its own MCP server, and a central host discovers and routes tool calls without hardcoding provider-specific logic. On top of this data layer, a LangGraph multi-agent system coordinates five specialist agents (recovery, training load, environmental context, route planning, and fitness knowledge) over the Agent-to-Agent (A2A) protocol. Each specialist is a Reasoning-and-Acting (ReAct) agent scoped to its own MCP servers; an orchestrator decomposes a query, delegates to the relevant specialists in parallel, and synthesises the answer. The fitness specialist uses Retrieval-Augmented Generation (RAG) over a vector database of public-domain exercise-science books, so its advice stays tied to a citable source. We evaluate the system end to end with MLflow's conversation-simulation framework, using ten synthetic personas across two user archetypes that between them exercise all five domains. Each conversation is scored on five dimensions: goal completion, user frustration, safety, supportive tone, and data grounding. The last is deterministic, checking tool use from the execution trace rather than the chat text. Adding a new data source takes one file and one configuration line, and the multi-agent design produces context-aware recommendations grounded in the athlete's own recovery, load, weather, and route data.